\section{All rational points of the one-fifth source}
\label{sec:one-fifth-points}

We now determine the rational points of the curves $C,D$ in
\eqref{eq:theta-C} and \eqref{eq:theta-quartic}.  The proof is an
explicit Dem'janenko--Manin comparison of two maps to one elliptic
curve; see also \cite{stoll2011}.  One fixed elliptic group is
certified by rational two-descent and saturation.  The height
comparison and the resulting finite endgame are given explicitly.

\begin{lemma}[The elliptic group]
\label{lem:fifth-group}
For
\[
 E:\quad y^2=x(x-81)(x-121),\qquad G=(49,336),
\]
one has
\[
 E(\Q)=\Z G\oplus E[2](\Q),\qquad
 E[2](\Q)=\{O,(0,0),(81,0),(121,0)\}.
\]
\end{lemma}

\begin{proof}
A two-descent gives a two-Selmer group with invariants $(2,2,2)$;
the torsion invariants are $(2,2)$ and $G$ has infinite order.
Thus the rank is one.  Full saturation gives
$H=(484/9,-8470/27)$, with $H+(121,0)=G$.
Here is the complete two-descent and saturation certificate, in Magma syntax:
\begin{verbatim}
E := EllipticCurve([Rationals() | 0,-202,0,9801,0]);
G := E![49,336,1];
T, tmap := TorsionSubgroup(E);
S, smap := TwoSelmerGroup(E);
assert Invariants(T) eq [2,2];
assert Invariants(S) eq [2,2,2];
assert Order(G) eq 0;
H, proved := Saturation([G] : TorsionFree := true);
assert proved and #H eq 1;
assert exists{t : t in [P @ tmap : P in T] |
             H[1] eq G+t or H[1] eq -G+t};
\end{verbatim}
The local-condition dimensions, starting with the supported
three-coordinate squareclass space, are $15,10,9,7,5,3,3$:
the successive steps are the norm condition and the places
$\infty,3,5,11,2$.  The prime-saturation bound is $7$; the
additional small-height search bound in the full saturation routine is $1379$.
This certificate uses the rational descent and full saturation
algorithms, with no analytic rank or GRH input.
\end{proof}

\begin{theorem}[The complete rational point sets]
\label{thm:fifth-complete-points}
The quartic and its genus-five double cover have rational point sets
\[
 D(\Q)=\{[\epsilon:\eta:1]:\epsilon,\eta\in\{1,-1\}\},
\]
\[
 C(\Q)=\{[\epsilon_1:\epsilon_2:\epsilon_3:1:0]:
                   \epsilon_i\in\{1,-1\}\}.
\]
Consequently no strict rational-distance point has an opposite-pair
distance difference equal to $1/5$ or $-1/5$.
\end{theorem}

\begin{proof}
There is no real point of $D$ with $Z=0$, since the homogeneous
quartic there is positive for $(X,Y)\ne(0,0)$.
Set $\xi=X/Z$, $\eta=Y/Z$ and $c=\xi^2+\eta^2-2$.
When $c=0$, the quartic gives $(\xi^2-\eta^2)^2=0$, so
$(\xi,\eta)=(\pm1,\pm1)$.
Suppose henceforth that $c\ne0$, and put
\[
 z=\frac{\xi^2-\eta^2}{14c},\qquad
 w_+=\frac{2\xi}{c},\qquad w_-=\frac{2\eta}{c}.
\]
Substitution gives
\begin{equation}
 w_+^2=(z^2+50)(z^2+14z+51),\qquad
 w_-^2=(z^2+50)(z^2-14z+51).
 \label{eq:fifth-two-maps}
\end{equation}
The first quartic is birational to $E$ through
\[
 A=w+z^2+7z+26,\qquad
 x=2A+49,\qquad y=(4z+14)A+336.
\]
Its inverse coordinate is
\begin{equation}
 \pi(P)=\frac{y(P)-7x(P)+7}{2(x(P)-49)}.
 \label{eq:fifth-pi}
\end{equation}
This function has simple poles at $O,-G$ and
$\pi(G)=-937/168$ at its removable singularity.
The involution over the $z$-line is $P\mapsto-G-P$.
Thus \eqref{eq:fifth-two-maps} produces rational points $P_+,P_-$
with $\pi(P_+)=z$ and $\pi(P_-)=-z$.

The map $P\mapsto3P+G$ commutes with this involution, and hence
descends to a degree-nine map $R$ on the $z$-line.  Write
$R=\mathcal A/\mathcal B$ with coefficients listed in
Table~\ref{tab:fifth-map}.  Substitution in the elliptic addition
law verifies
\begin{equation}
 R(\pi(P))=\pi(3P+G).
 \label{eq:fifth-triplication}
\end{equation}
Let $h_2([a:b])=\log_2\max(|a|,|b|)$ for coprime integers
$a,b$.  The B\'ezout coefficients displayed below give
\[
 U_0\mathcal A+V_0\mathcal B=D_0,
 \qquad U_\infty\mathcal A^\vee+V_\infty\mathcal B^\vee=D_0,
 \qquad D_0=55399054482258080041175062500,
\]
where $F^\vee(z)=z^9F(1/z)$ and all four B\'ezout polynomials
have degree at most eight.  Their paired coefficient norms, and
the norms of $\mathcal A,\mathcal B$, are less than $2^{124}$.
The two polynomials have no common root.

After homogenization, these identities give $D_0b^{17}$ and
$D_0a^{17}$.  The gcd of
$\mathcal A_h(a,b),\mathcal B_h(a,b)$ therefore divides $D_0$.
Use the identity corresponding to the larger of $|a|,|b|$ to
obtain the lower bound; the coefficient norms give the upper bound.
It follows that
\begin{equation}
 |h_2(R(t))-9h_2(t)|\leq124\qquad(t\in\PP^1(\Q)).
 \label{eq:fifth-height-error}
\end{equation}
The telescoping limit
$\widehat h_R(t)=\lim_{j\to\infty}9^{-j}h_2(R^j(t))$
therefore exists and differs from $h_2(t)$ by at most $124/8$.
The first four iterates of $\infty$ have maximum-coordinate
bit lengths $10,95,860,7754$.  Consequently
\begin{equation}
 h_0:=\widehat h_R(\infty)
 \geq\frac{7753-124/8}{9^4}
 =\frac{15475}{13122}>0.
 \label{eq:fifth-positive-height}
\end{equation}

Choose $A_0$ over a finite extension with $2A_0=G$.
The pole divisor $O+(-G)$ of $\pi$ is linearly equivalent to
$2(-A_0)$.  The standard comparison with the canonical elliptic
height, together with
$(3P+G)+A_0=3(P+A_0)$, gives
\[
 \widehat h_R(\pi(P))=2\widehat h_E(P+A_0).
\]
Here the heights use logarithms to base two; the bounded error in
the divisor-height comparison disappears in the limit.
Lemma~\ref{lem:fifth-group} now yields
\begin{equation}
 \widehat h_R(\pi(mG+T))=(2m+1)^2h_0
 \quad(m\in\Z,\ T\in E[2](\Q)).
 \label{eq:fifth-height-multiple}
\end{equation}
Write $P_+=mG+T$ and $P_-=nG+T'$.
Since $h_2(z)=h_2(-z)$, the height error gives
\[
 4|(m-n)(m+n+1)|h_0\leq31,
 \qquad |(m-n)(m+n+1)|<7.
\]

If the product is zero, the involution $P\mapsto-G-P$
reduces to $n=m$, and a two-torsion translation must send $z$
to $-z$.  The complete translation table is
\[
\begin{array}{c|c|c}
T''&\pi(P+T'')&\text{equation for finite }z\\\hline
O&z&z=0\\
(0,0)&(-z+700)/(14z+1)&z^2=-50\\
(81,0)&(-35z-205)/(4z+35)&z^2=205/4\\
(121,0)&(-35z-305)/(6z+35)&z^2=305/6.
\end{array}
\]
The last two equations fail at $5$, the second fails over $\R$,
and $z=0$ would give $w_+^2=2550$, which fails at $2$.

If the product is nonzero, write $a=m-n$, $b=m+n+1$.
Then $|a|+|b|\leq|ab|+1\leq7$, whence $-4\leq m\leq3$.
The involution replaces $m$ by $-m-1$ without changing $z$,
so it remains to inspect $P_+=mG+T$ with $0\leq m\leq3$.
For $z=A/B$ in lowest terms the second equation of
\eqref{eq:fifth-two-maps} requires
\[
 \mathcal N=(A^2+50B^2)(A^2-14AB+51B^2)
\]
to be an integer square.  The following table, obtained by the
rational addition law on $E$, gives $(p,v_p(\mathcal N),u)$,
where $u$ is the nonzero unit residue modulo $p$.
Every displayed $u$ is a quadratic nonresidue and every valuation
is even.  The missing entry is $z=\infty$.
\[
\begin{array}{c|cccc}
 m&T=O&T=(0,0)&T=(81,0)&T=(121,0)\\\hline
 0&\infty&(3,4,2)&(3,4,2)&(11,2,7)\\
 1&(5,0,3)&(3,4,2)&(3,4,2)&(11,2,2)\\
 2&(5,0,3)&(3,4,2)&(3,4,2)&(11,2,8)\\
 3&(11,0,2)&(3,4,2)&(3,4,2)&(11,2,8).
\end{array}
\]
This excludes every finite $z$, proving the assertion about $D$.

The quotient $C\to D$ is $[P:Q:R:z:T]\mapsto[Q:R:P]$.
Above the four points just found, $Q^2=R^2=P^2$.
The equations of $C$ give $28zT=0$ and $14zT+T^2=0$,
so $T=0$.  The eight stated boundary points are exactly these
fibres.  The strict return \eqref{eq:theta-return} uses $T=1$,
which is impossible.  The square symmetries give both opposite
pairs and both signs of the difference.
\end{proof}

\begin{center}
\small
\refstepcounter{table}\label{tab:fifth-map}
\begin{tabular}{r|r r}
$i$ & $[z^i]\mathcal A$ & $[z^i]\mathcal B$\\\hline
0 & 111607125000 & -12021749375\\
1 & 82021123125 & -6423322500\\
2 & 26108197500 & -825405000\\
3 & 5174999950 & 193235350\\
4 & 768860400 & 85090800\\
5 & 79894350 & 15797250\\
6 & 3064600 & 2089798\\
7 & -389550 & 206808\\
8 & -52500 & 12678\\
9 & -1874 & 336\\
\end{tabular}
\par\smallskip
Table~\thetable: The degree-nine map in \eqref{eq:fifth-triplication}.
\end{center}
\noindent The following two tables list the Bezout coefficients in increasing degree.
Each row gives the coefficients of the two indicated polynomials.
\begin{center}
\scriptsize
\begin{tabular}{r|r r}
$i$ & $[z^i]U_0$ & $[z^i]V_0$\\\hline
0 & -2997498492457923614300 & -27832687036826975619900\\
1 & -1095858706252693720500 & -15753573106648371231300\\
2 & -19762427635601063700 & -3841783019758516568700\\
3 & 51931145761277052296 & -636028100318545578700\\
4 & 12434000961572564136 & -83226348410177026200\\
5 & 1819310074540097136 & -5702224020804565800\\
6 & 210502428872007552 & 220117172270658448\\
7 & 15563169683956800 & 60607066424472960\\
8 & 483327205240320 & 2695699948274880\\
\end{tabular}
\end{center}
\begin{center}
\scriptsize
\begin{tabular}{r|r r}
$i$ & $[z^i]U_\infty$ & $[z^i]V_\infty$\\\hline
0 & -190236680175253051840321031250 & -1060857558315303455561667968750\\
1 & -5532338056063458586310158593750 & -20552009231971099412950192968750\\
2 & -69059947929901090666818928125000 & -41727935009628484047973450781250\\
3 & -580126053198122963653101514843750 & 2117256664470905046006824705468750\\
4 & -3869791759426723943158100156250000 & 26897617537416601772185292900390625\\
5 & -14274207303033339557638111728515625 & 200210932783826079838895782705078125\\
6 & 16848129423821193386849827861328125 & 1177195891501011517791195006591796875\\
7 & 326540805971266624719922115478515625 & 4456602395261102375145481212158203125\\
8 & 765209073668489627429760501708984375 & 7104023056216253086689948486328125000\\
\end{tabular}
\end{center}

